\chapter{General Equilibrium Theory: Pure Exchange Economy}

\section{Introduction to General Equilibrium}

Economic equilibrium can be analyzed using two main approaches.

\subsection{Partial Equilibrium}
Partial equilibrium analysis focuses on a single market while holding all other factors constant (ceteris paribus).  
This approach provides simple and intuitive insights.

\subsection{General Equilibrium}
General equilibrium analysis studies all markets simultaneously, removing the ceteris paribus assumption.  
It captures interdependence across markets and provides a more complete framework.

\paragraph{Definition}
A \textbf{general equilibrium} is a situation in which:
\begin{itemize}
    \item Consumers maximize utility subject to budget constraints
    \item All markets clear (demand equals supply)
    \item All agents act optimally given prices
\end{itemize}

\paragraph{Intuition}
General equilibrium can be interpreted as a state in which a single price system
coordinates all decisions in the economy.  
Consumers, taking prices as given, choose optimal bundles, and these choices are
mutually consistent across all markets.

---

\section{Walrasian Equilibrium}

A \textbf{Walrasian equilibrium} is a set of prices and allocations such that:
\begin{itemize}
    \item Consumers maximize utility
    \item Markets clear simultaneously
\end{itemize}

\subsection{Walras' Law}

Let $Z_i(p)$ denote excess demand for good $i$. Then:
\[
\sum_{i=1}^{n} p_i Z_i(p) = 0
\]

\paragraph{Intuition}
Walras' Law states that the total value of excess demand in the economy is always zero.  
If there is excess demand in one market, there must be excess supply in another.

\paragraph{Implication}
If $(n-1)$ markets are in equilibrium, the last market must also be in equilibrium.
Thus, only $(n-1)$ independent equilibrium conditions are needed.

---

\section{Pure Exchange Economy}

We consider an economy with:
\begin{itemize}
    \item Two consumers: $A, B$
    \item Two goods: $X, Y$
\end{itemize}

\subsection{Initial Endowment}

\[
\omega_A = (X_A^0, Y_A^0), \quad
\omega_B = (X_B^0, Y_B^0)
\]

\[
\bar{X} = X_A^0 + X_B^0, \quad
\bar{Y} = Y_A^0 + Y_B^0
\]

\paragraph{Interpretation}
The total amount of goods is fixed, and the problem is how to allocate these goods
between consumers.

---

\subsection{Preferences}

\[
U_A = U_A(X_A, Y_A), \quad
U_B = U_B(X_B, Y_B)
\]

Preferences are assumed to be continuous and strictly convex.

\paragraph{Interpretation}
Convex preferences imply diminishing marginal rate of substitution, ensuring
interior solutions and smooth trade-offs between goods.

---

\section{Edgeworth Box}

The Edgeworth box represents all feasible allocations:
\[
X_A + X_B = \bar{X}, \quad
Y_A + Y_B = \bar{Y}
\]

Each point corresponds to an allocation between the two consumers.

\subsection{Initial Allocation}

The initial endowment point is:
\[
I = (\omega_A, \omega_B)
\]

\paragraph{Intuition}
The Edgeworth box illustrates all possible reallocations of goods.  
Starting from the initial endowment, trade moves the economy to other feasible allocations.

---

\section{Pareto Efficiency and Contract Curve}

\subsection{Pareto Efficiency}

An allocation is Pareto efficient if no individual can be made better off without making someone else worse off.

\paragraph{Intuition}
At a Pareto efficient allocation, all mutually beneficial trades have been exhausted.

---

\subsection{Contract Curve}

The contract curve is the set of Pareto efficient allocations:
\[
MRS_A = MRS_B
\]

where
\[
MRS_i = \frac{\partial U_i / \partial X_i}{\partial U_i / \partial Y_i}
\]

\paragraph{Intuition}
If $MRS_A \neq MRS_B$, then gains from trade still exist.  
Equality implies no further mutually beneficial exchange is possible.

---

\section{Exchange and Price Adjustment}

A Walrasian auctioneer announces prices $p = (p_X, p_Y)$.

Each consumer solves:
\[
\max U_i(X_i, Y_i)
\quad \text{s.t.} \quad
p_X X_i + p_Y Y_i = p_X X_i^0 + p_Y Y_i^0
\]

\paragraph{Interpretation}
Consumers choose bundles that maximize utility given their income (value of endowment).

---

\subsection{Excess Demand}

\[
Z_X = (X_A + X_B) - \bar{X}
\]
\[
Z_Y = (Y_A + Y_B) - \bar{Y}
\]

Prices adjust until:
\[
Z_X = 0, \quad Z_Y = 0
\]

\paragraph{Intuition}
If demand exceeds supply, prices rise; if supply exceeds demand, prices fall.  
This adjustment process continues until markets clear.

---

\section{Consumer Optimization}

\subsection{Agent A}

\[
\max U_A(X_A, Y_A)
\]
\[
\text{s.t. } p_X X_A + p_Y Y_A = p_X X_A^0 + p_Y Y_A^0
\]

\[
\mathcal{L}_A =
U_A(X_A, Y_A)
+ \lambda_A (p_X X_A^0 + p_Y Y_A^0 - p_X X_A - p_Y Y_A)
\]

\[
\frac{\partial U_A}{\partial X_A} = \lambda_A p_X
\quad
\frac{\partial U_A}{\partial Y_A} = \lambda_A p_Y
\]

\[
MRS_A = \frac{p_X}{p_Y}
\]

\paragraph{Interpretation}
Consumers equate their marginal rate of substitution to the price ratio.  
This ensures that the subjective trade-off equals the market trade-off.

---

\subsection{Agent B}

\[
MRS_B = \frac{p_X}{p_Y}
\]

---

\section{General Equilibrium Condition}

\[
MRS_A = MRS_B = \frac{p_X}{p_Y}
\]

\[
Z_X = 0, \quad Z_Y = 0
\]

\paragraph{Key Insight}
General equilibrium simultaneously achieves:
\begin{itemize}
    \item Individual optimality
    \item Market clearing
    \item Pareto efficiency
\end{itemize}

---

\section{Relative Prices}

Only relative prices matter:
\[
\frac{p_X}{p_Y}
\]

We may normalize:
\[
p_X + p_Y = 1
\]

\paragraph{Interpretation}
Absolute prices are irrelevant; only price ratios affect decisions.

---

\section{Existence of Equilibrium}

Equilibrium exists if:
\begin{itemize}
    \item Excess demand functions are continuous
    \item Goods are desirable
\end{itemize}

Then there exists $p^*$ such that:
\[
Z(p^*) = 0
\]

---

\section{Walras' Law in Two-Good Case}

\[
p_X Z_X + p_Y Z_Y = 0
\]

\paragraph{Implication}
If one market clears, the other must also clear automatically.

---

\section{Economic Interpretation}

General equilibrium implies:
\begin{itemize}
    \item All markets clear simultaneously
    \item Agents optimize given prices
    \item Allocation is Pareto efficient
\end{itemize}

\paragraph{Final Insight}
Prices act as a coordination mechanism that aligns individual incentives with
aggregate feasibility, leading to an efficient allocation of resources.